\subsection{Hoja de datos}
\label{subsec:hoja-de-datos}

Se adjunta la hoja de datos original con las mediciones realizadas en el laboratorio.

\subsubsection{Práctica 1 --- Ley de Ohm}

\begin{table}[H]
    \centering
    \caption{Medición de corriente y voltaje con conectores en los extremos}
    \label{tab:resistencia_1}
    \begin{tabular}{|c|c|}
        \hline
        \textbf{(disp A)I (A)} & \textbf{(disp B)V (V)} \\ \hline
        \csvreader[
            separator=semicolon,
            head=true,
            late after line=\\ \hline
        ]{secciones/hoja-datos/resistencia_1.csv}{}
        {\csvcoli & \csvcolii}
    \end{tabular}
\end{table}

\begin{table}[H]
    \centering
    \caption{Medición de corriente y voltaje con conector en el centro}
    \label{tab:resistencia_2}
    \begin{tabular}{|c|c|}
        \hline
        \textbf{(disp A)I (A)} & \textbf{(disp B)V (V)} \\ \hline
        \csvreader[
            separator=semicolon,
            head=true,
            late after line=\\ \hline
        ]{secciones/hoja-datos/resistencia_2.csv}{}
        {\csvcoli & \csvcolii}
    \end{tabular}
\end{table}

\subsubsection{Práctica 2 --- Medición de Resistencias}

\begin{table}[H]
    \centering
    \caption{Mediciones directas con multímetro UT58D (modo $\Omega$)}
    \label{tab:hoja-resistencias}
    \begin{tabular}{|c|c|c|c|}
        \hline
        \textbf{Resistencia} & \textbf{Rango utilizado} & \textbf{$R_{med}$} & \textbf{$\Delta R$} \\ \hline
        R1 & 20 M$\Omega$   & 5{,}61 M$\Omega$  & 0{,}1 M$\Omega$ \\ \hline
        R2 & 2000 k$\Omega$ & 562 k$\Omega$     & 8 k$\Omega$ \\ \hline
        R3 & 20 k$\Omega$   & 5{,}50 k$\Omega$  & 0{,}05 k$\Omega$ \\ \hline
    \end{tabular}
\end{table}

\subsubsection{Práctica 3 --- Influencia del aparato de medida}
 
\begin{table}[H]
    \centering
    \caption{Mediciones de tensión con multímetro UT58D (modo V) --- Práctica 3}
    \label{tab:hoja-p3}
    \begin{tabular}{|c|c|c|c|c|c|}
        \hline
        \textbf{Circ.} & $R_a$ & $R_b$
            & \textbf{Rango} & $V_{Ra}$ & $V_{Rb}$ \\ \hline
        A & 5{,}6 k$\Omega$ & 5{,}6 M$\Omega$
            & 200\,mV / 20\,V & 5{,}0\,mV & 5{,}12\,V \\ \hline
        B & 560 k$\Omega$ & 5{,}6 M$\Omega$
            & 200\,mV / 200\,mV & 0{,}20\,mV & 2{,}5\,mV \\ \hline
    \end{tabular}
\end{table}